% ===== Block B08: 原文 p.67–72 | §3.4 关联函数 | 译手 tr-ch3b =====

\subsection{关联函数}

% [原文 p.67]
设 $C$ 为半稳定曲线（semi-stable curve），$\omega_C$ 为其对偶化层（dualizing sheaf）。置
$C^M=\underbrace{C\times\cdots\times C}_{M}$。则 $C^M$ 具有余维 $1$ 的奇点，但由于 $C^M$
局部是完全交（complete intersection），我们仍可定义对偶化层 $\omega_{C^M}$。
（例如见 [BS1] 或 [Kl]。）此外，我们可以证明
\[
  \omega_{C^M} = \omega_C^{\boxtimes M}
\]
其中 $\pi_j : C^M \to C$ 是第 $j$ 个投影，且我们定义
\[
  \omega_C^{\boxtimes M} = \pi_1^*\omega_C \otimes \pi_2^*\omega_C
  \otimes \cdots \otimes \pi_M^*\omega_C .
\]
% [Tutor Note: $\boxtimes$ 是「外张量积」（exterior tensor product），与通常的
%  $\otimes$（在同一空间上取张量积）不同：这里各 $\pi_j^*\omega_C$ 分别来自 $C^M$ 的不同
%  因子，故必须用 $\boxtimes$。核对方法：取 $M=1$ 时 $\omega_C^{\boxtimes 1}=\pi_1^*\omega_C$，
%  与 $\omega_{C^1}=\omega_C$ 一致；取 $C$ 光滑、$M=2$ 时 $\omega_{C\times C}=
%  \mathrm{pr}_1^*\omega_C\otimes\mathrm{pr}_2^*\omega_C$ 即 Künneth 公式的二维特例。原文 p.67]
%
% [Tutor Note: 「$C^M$ 局部是完全交」是该页成立的关键：完全交的概形上对偶化层
%  总可用 Koszul 复形定义（$\omega_{C^M}=\omega_{\PP^N}|_{C^M}$ 的局部版本），所以
%  $C^M$ 虽有奇点（两条分支在结点的乘积处相交，余维 $1$），$\omega_{C^M}$ 仍存在。
%  注意这里奇点是「余维 $1$」而非余维 $\ge 2$，故不能靠 Hartogs 型延拓绕过。原文 p.67]
（例如见 [Kl]。）由于对奇异的半稳定曲线 $C$，$C^M$ 有奇点，$C^M$ 的第 $(i,j)$ 条对角线
$\Delta_{ij}=\{(P_1,\dots,P_N)\mid P_i=P_j\}$ 只是 Weil 除子（Weil divisor），
而不是 Cartier 除子（Cartier divisor）。设
$\mathfrak{X}=(C;Q_1,\dots,Q_N;\eta_1,\dots,\eta_N)$ 是带形式邻域的稳定 $N$-点曲线。
% [Tutor Note: 这里 $\Delta_{ij}$ 的记法照原文排，但 $C^M$ 的点的写法应为
%  $(P_1,\dots,P_M)$（$M$ 个分量），原文两处（本页及定理 3.24 中）都印成
%  $(P_1,\dots,P_N)$，见回报 ERRATA。判断依据：$N$ 是 $C$ 上的标点数，与 $C^M$
%  的分量个数无关。原文 p.67]
%
% [Tutor Note: 为什么是「Weil 除子而非 Cartier 除子」：$\Delta_{ij}\subset C^M$ 在
%  $C$ 有结点（$C$ 奇异的半稳定曲线）时自身含有 $C^M$ 的奇点，故它不是
%  Cartier 除子，不能直接写 $\mathcal{O}(\Delta_{ij})$；于是定理 3.24 里用
%  $*\Delta_{ij}$（星号 = 取该除子的「真正原像/严格变换」型拉回，见定理 3.24 的
%  $H^0$ 表达式），这样 $\omega^{\boxtimes M}(*\Delta_{ij}+\cdots)$ 才是良定义的层。
%  自检：$C$ 光滑时 $\Delta_{ij}$ 是光滑除子，$*$ 可去掉。原文 p.67]

\begin{theorem}[3.24]
固定 $\langle\Psi|\in\cV_\lambda^\dagger(\mathfrak{X})$。对每个非负整数 $M$，数据
\[
  X_1,X_2,\dots,X_M\in\mathfrak{g},\quad |\Phi\rangle\in\cH_{\mathfrak{X}}
\]
确定一个元素
\[
  F = \langle\Psi|X_1(P_1)X_2(P_2)\cdots X_M(P_M)|\Phi\rangle
      dP_1dP_2\cdots dP_M
\]
属于
\[
  H^0\Bigl(C^M,\ \omega_{C^M}^{\boxtimes M}\Bigl(
    \sum_{1\le i<j\le M}*\Delta_{ij}
    +\sum_{i=1}^{M}\sum_{j=1}^{N}*\pi_i^{-1}(Q_j)\Bigr)\Bigr),
\]
其中 $\Delta_{ij}=\{(P_1,\dots,P_N)\mid P_i=P_j\}$ 是对角线（diagonal）。这个亚纯形式
具有如下性质。

(0)\quad 当 $M=0$ 时，$F=\langle\Psi|\Phi\rangle$ 是由配对 (3.12) 诱导的典范配对。

(1)\quad $F$ 关于 $|\Phi\rangle$ 是线性的，关于 $X_i$ 是多线性的。

(2)\quad 对任意置换 $\sigma\in\mathfrak{S}_M$，我们有
\[
  F = \langle\Psi|X_{\sigma(1)}(P_{\sigma(1)})X_{\sigma(2)}(P_{\sigma(2)})
      \cdots X_{\sigma(M)}(P_{\sigma(M)})|\Phi\rangle dP_1dP_2\cdots dP_M .
\]
例如，对换位 $(i,i+1)$，我们有
\[
  \begin{aligned}
  F &= \langle\Psi|X_1(P_1)\cdots X_{i-1}(P_{i-1})X_{i+1}(P_{i+1})X_i(P_i)\\
    &\quad X_{i+2}(P_{i+2})\cdots X_M(P_M)|\Phi\rangle dP_1dP_2\cdots dP_M .
  \end{aligned}
\]

(3)\quad 对 $k=1,\dots,N$ 及 $\xi_k=\eta_k^{-1}(\xi)$，若 $\xi_k$ 是全纯坐标，
则我们有等式
\[
\begin{aligned}
&\operatorname{Res}_{\xi_k=0}\bigl(\xi_k^n\langle\Psi|X(\xi_k)X_1(P_1)X_2(P_2)
  \cdots X_M(P_M)|\Phi\rangle d\xi_k\bigr)\\
&\qquad= \langle\Psi|X_1(P_1)X_2(P_2)\cdots X_M(P_M)|\rho_k(X(n))\Phi\rangle .
\end{aligned}
\]
% [Tutor Note: 性质 (3) 是「$T(z)$ 与 $X(z)$ 的算子积展开的留数形式」：左端把
%  $\xi_k$ 方向的第 $n$ 个洛朗系数取出，右端就是 $\rho_k(X(n))$ 作用在 $|\Phi\rangle$ 上。
%  符号核对：$\xi_k^{-n-1}d\xi_k$ 的留数 = 第 $n$ 个系数，故 $\xi_k^n$ 配 $X(n)$；
%  若把 $\xi_k^n$ 写成 $\xi_k^{n+1}$ 就会偏移一位（下一页定理 3.26 的 $L_n$ 正是
%  $\xi_k^{n+1}$，因为 $d\xi_k^2$ 的留数偏移一格）。原文 p.67]
换言之，我们有展开式
\[
  \begin{aligned}
  &\langle\Psi|X(\xi_k)_1(P_1)X_2(P_2)\cdots X_M(P_M)|\Phi\rangle d\xi_k\\
  &\quad = \sum_{n\in\Z}\langle\Psi|X_1(P_1)X_2(P_2)\cdots X_M(P_M)
     |\rho_k(X(n))\Phi\rangle\xi_k^{-n-1}d\xi_k .
  \end{aligned}
\]
% [Tutor Note: 上式左端的 $X(\xi_k)_1(P_1)$ 照原文排；按前后文（性质 (3) 与
%  定理 3.26 的同型展开）它应为 $X(\xi_k)X_1(P_1)$，原文漏了一个 $X_1$ 的记号，
%  见回报 ERRATA。原文 p.67]
% [Tutor Note: 定理 3.24 中 $H^0$ 里的极点除子有两类：$*\Delta_{ij}$ 给出
%  $X_i,X_j$ 两个插入点重合时的二阶极点（这正是性质 (4) 的 $(z-z')^{-2}$ 项），
%  $*\pi_i^{-1}(Q_j)$ 给出第 $i$ 个插入点趋近标点 $Q_j$ 时的极点（性质 (5)）。
%  维数检验：$M$ 个点各贡献一个 $dP_i$，共 $M$ 个一次微分；而 $\omega_{C^M}^{\boxtimes M}$
%  的截面是 $\bigotimes_i \omega_{C,P_i}$，故 $F$ 是 $C^M$ 上的全纯 $M$-形式乘以极点，
%  这与 $M=0$ 时退化为常数 $\langle\Psi|\Phi\rangle$ 相容。原文 p.67]

% [原文 p.68]
(4)\quad 对曲线 $C$ 的非奇异点 $P$ 处的局部全纯坐标 $z$，我们有如下等式。
\[
  \begin{aligned}
  &\langle\Psi|X(P)Y(P')X_1(P_1)X_2(P_2)\cdots X_M(P_M)|\Phi\rangle\\
  &\quad = \frac{\ell\cdot(X,Y)}{(z(P)-z(P'))^2}
      \langle\Psi|X_1(P_1)X_2(P_2)\cdots X_M(P_M)|\Phi\rangle\\
  &\quad\quad +\frac{1}{z(P)-z(P')}
      \langle\Psi|[X,Y](P')X_1(P_1)X_2(P_2)\cdots X_M(P_M)|\Phi\rangle\\
  &\quad\quad +\text{在 } P=P' \text{ 处正则}.
  \end{aligned}
\]
% [Tutor Note: 这里 $\ell\cdot(X,Y)$ 是原书对 $\mathfrak{g}$ 上不变双线性型的记法：
%  $\ell$ 是仿射李代数的水平（level），$(\cdot,\cdot)$ 是归一化不变型。核对方法：
%  取 $\mathfrak{g}=\mathfrak{sl}_2$、$X=Y=J^3$，则标准 OPE 为
%  $J^3(z)J^3(w)\sim \ell/2\,(z-w)^{-2}$，与 $\ell\cdot(X,Y)$ 在
%  $(J^3,J^3)=1/2$ 时一致。
%  另：性质 (4) 的符号是「物理约定」：$T$ 与 $X$ 的 OPE 全纯部分给出
%  $[X,Y]$ 项。自检办法：令 $X=Y$，则 $[X,X]=0$，等式退化为
%  $\langle X(P)X(P')\rangle=\ell\cdot(X,X)/(z-z')^2+\text{正则}$，
%  这正是自由玻色子的两点函数。原文 p.68]

(5)\quad 对 $Q_i$ 处的局部全纯坐标 $z$，以及对
$|v\rangle\in V_{\mathfrak{X}}^\dagger=V_{\lambda_1}\otimes\cdots\otimes V_{\lambda_N}
\subset\cH_{\mathfrak{X}}$，我们有等式
\[
  \begin{aligned}
  &\langle\Psi|X(P)X_1(P_1)X_2(P_2)\cdots X_M(P_M)|v\rangle\\
  &\quad = \frac{1}{z(P)-z(Q_i)}
      \langle\Psi|X_1(P_1)X_2(P_2)\cdots X_M(P_M)|\rho_i(X)v\rangle\\
  &\quad\quad +\text{在 } P=Q_i \text{ 处正则}.
  \end{aligned}
\]
% [Tutor Note: (5) 中 $|v\rangle$ 只取有限维张量积
%  $V_{\lambda_1}\otimes\cdots\otimes V_{\lambda_N}$ 的元素，而不是整个
%  $\cH_{\mathfrak{X}}$：因为只有对最高权模 $V_{\lambda_i}$ 里的向量，
%  $\rho_i(X)|v\rangle$ 才有意义（$\cH_{\mathfrak{X}}$ 含 $\mathfrak{g}$ 的
%  无穷维部分 $X(-1)|0\rangle$ 之类，$\rho_i$ 不作用其上）。这一点在下一页
%  Claim 1 里用到：那里 $|\tilde u\rangle$ 可以含 $X(-1)|0\rangle$。
%  原文 p.68]

\end{theorem}
这些函数 $F$ 称为流的关联函数（correlation functions of currents）。

\begin{sloppypar}\noindent\textit{证明.}\quad 选取曲线 $C$ 的 $M+1$ 个非奇异点 $P_1,P_2,\dots,P_M,P$
及其形式邻域 $\eta_{N+1},\eta_{N+2},\dots,\eta_{N+M+1}$。置
\end{sloppypar}
\[
  \hat{\mathfrak{X}} = (C;Q_1,\dots,Q_N,Q_{N+1},\dots,Q_{N+M+1};
      \eta_1,\dots,\eta_{N+M+1})
\]
\[
  \tilde{\mathfrak{X}} = (C;Q_1,\dots,Q_N,Q_{N+1},\dots,Q_{N+M};
      \eta_1,\dots,\eta_{N+M})
\]
其中 $Q_{N+i}=P_i$，$i=1,\dots,M$，且 $Q_{N+M+1}=P$。由定理 3.24，存在典范同构
% [Tutor Note: 原文此处写「由定理 3.24」，但这一段正是定理 3.24 的证明，
%  不能引用自身；所指应为 §3.3 中「增加标号 0 的标点后共形块空间同构」的
%  定理 3.15（多点情形见推论 3.16；3.23 是命题，与本同构无关）。见报告 ERRATA。原文 p.68]
\[
  \iota_M : \cV_\lambda^\dagger(\tilde{\mathfrak{X}})
  \simeq \cV_{\lambda,\vec{0}_M}^\dagger(\hat{\mathfrak{X}})
\]
\[
  \iota_{M+1} : \cV_{\lambda,\vec{0}_M}^\dagger(\tilde{\mathfrak{X}})
  \simeq \cV_{\lambda,\vec{0}_{M+1}}^\dagger(\hat{\mathfrak{X}})
\]
其中 $\vec{0}_k=(\underbrace{0,\dots,0}_{k})$。对
$\langle\Psi|\in\cV_\lambda^\dagger(\tilde{\mathfrak{X}})$，置
\[
  \langle\tilde{\Psi}| = \iota_M(\langle\Psi|),\qquad
  \langle\hat{\Psi}| = \iota_{M+1}(\langle\tilde{\Psi}|).
\]
% [Tutor Note: 两个同构的分工要点：$\iota_M$ 把「$\tilde{\mathfrak{X}}$ 上的
%  $\Psi$」搬到「$\hat{\mathfrak{X}}$ 上、且在 $P$ 处取平凡权 $0$」；
%  $\iota_{M+1}$ 再把它搬到「在 $P'$ 处也取平凡权 $0$」。串联起来
%  $\langle\hat{\Psi}|=\iota_{M+1}\iota_M\langle\Psi|$ 才是「在 $P,P'$ 两点都
%  插入平凡表示」的真空态，故可用来算 $X(-1)Y(-1)$ 的矩阵元。
%  自检：$\vec{0}_0=()$，$M=0$ 时 $\iota_0$ 就是恒等同构。原文 p.68]
\hfill$\square$
% [Tutor Note: 定理 3.24 的证明在此处只给出了「把 $\langle\Psi|$ 搬到扩充的
%  曲线上」这一准备，真正的构造分三步，见下文 Claim 1--3；原文把三段论证
%  夹在证明中间，读时须注意 Claim 1--3 各自独立带证毕号，最后在 p.70 末尾
%  回到本证明。原文 p.68]

\noindent\textbf{断言 1.}\quad 对任意
$|\tilde{u}\rangle\in\cH_{\tilde{\mathfrak{X}}}\otimes\cH_{\vec{0}_M}$ 与
$X\in\mathfrak{g}$，$\langle\hat{\Psi}|\tilde{u}\otimes X(-1)|0\rangle d\eta$
定义了曲线 $C$ 在点 $P$ 处的余切向量。

\noindent\textit{证明.}\quad 选取 $C$ 上的亚纯函数
$f\in H^0(C,\cO_C(*(P+Q_1)))$ 使得
\[
  f = \eta^{-1} + \text{在 } P \text{ 处正则},
\]
\[
  f\equiv 0 \bmod(\xi_j^{n_j}) \quad \text{在 } Q_j \text{ 处},\ j\ne 1,
\]
其中 $\eta=\eta_{M+1}^{-1}(\xi)$，$\xi_j=\eta_j^{-1}(\xi)$，且 $n_j$ 充分大，
使得 $\rho_j(X[f])|\tilde{u}\rangle=0$ 且 $f$ 在 $Q_j$ 处全纯，$j\ne 1$。于是有
\[
  \langle\hat{\Psi}|\tilde{u}\otimes X(-1)|0\rangle
  = \langle\hat{\Psi}|\tilde{u}\otimes(X[f])|0\rangle
  = -\langle\hat{\Psi}|\rho_1(X[f])\tilde{u}\otimes 0\rangle .
\]
% [Tutor Note: 第一个等号是本段的关键：由 $f=\eta^{-1}+\text{正则}$ 得
%  $X[f]=X(-1)+\sum_{n\ge 0}c_nX(n)$，而 $X(n)|0\rangle=0$（$n\ge 0$），故
%  $X[f]|0\rangle=X(-1)|0\rangle$，即把 $X(-1)$ 换成「沿 $f$ 的抹平」后矩阵元不变。
%  第二个等号把 $X[f]$ 从左边移到右边，代价是 $\rho_1$：这正是真空的 Ward 恒等式
%  $\langle X[f]\cdots\rangle=-\langle\cdots\rho_1(X[f])\rangle$。注意
%  $f$ 的极点条件 $f\equiv0 \bmod \xi_j^{n_j}$ 是为了让 $\rho_j(X[f])|\tilde u\rangle=0$，
%  即让除 $Q_1$ 以外的贡献全部消失。原文 p.68]
%  另：这里的 $\eta_{M+1}$、$\tilde\eta_{M+1}$ 与上一段的
%  $\eta_{N+1},\dots,\eta_{N+M+1}$ 指标不一致：$P=Q_{N+M+1}$，故 $P$ 处的形式邻域
%  应为 $\eta_{N+M+1}$。原文可能把「新增加的第 $M+1$ 个点」简记为 $M+1$，
%  但按上一段的记法应写 $N+M+1$，见回报 ERRATA。原文 p.68]
因此，若把形式邻域 $\eta_{M+1}$ 换成 $\tilde{\eta}_{M+1}$，我们就有
\[
  \tilde{\eta} = \tilde{\eta}_{M+1}^{-1}(\xi)
  = a_1\eta + a_2\eta^2 + \cdots,\quad a_1\ne 0 .
\]

% [原文 p.69]
\[
  \langle\hat{\Psi}|\tilde{u}\otimes X(-1)|0\rangle_{\tilde{\eta}}
  = a_1^{-1}\langle\hat{\Psi}|\tilde{u}\otimes X(-1)|0\rangle_{\eta} .
\]
% [Tutor Note: 系数是 $a_1^{-1}$ 而不是 $a_1$：$d\tilde\eta=a_1d\eta+\cdots$，
%  而不变量是 $\langle\cdots\rangle_\eta d\eta$，故 $\langle\cdots\rangle$
%  按 $a_1^{-1}$ 缩放。这一条正是「$\langle\cdots\rangle_\eta d\eta$ 与形式坐标
%  选取无关」的实质，直接给出一阶无穷小邻域上的余切向量。核对：取
%  $\tilde\eta=2\eta$（$a_1=2$），则 $\langle\cdots\rangle_{\tilde\eta}
%  =\frac12\langle\cdots\rangle_\eta$，而 $d\tilde\eta=2d\eta$，乘积不变。
%  原文 p.69]
这蕴涵 $\langle\hat{\Psi}|\tilde{u}\otimes X(-1)|0\rangle$ 只依赖于一阶无穷小邻域，
并且 $\langle\hat{\Psi}|\tilde{u}\otimes X(-1)|0\rangle_\eta d\eta\in T_P^*C$
与形式坐标的选取无关。\hfill$\square$

\noindent\textbf{断言 2.}\quad 置
\[
  \omega_j = \sum_{n\in\Z}(\tilde{\Psi}|\rho_j(X(n))|\tilde{u}\rangle
  \xi_j^{-n-1}d\xi_j,\qquad j=1,2,\dots,N+M,
\]
其中 $\xi_j=\eta_j^{-1}(\xi)$。则 $C$ 上存在亚纯 $1$-形式
\[
  \omega\in H^0\Bigl(C,\ \omega_C\bigl(*\sum_{j=1}^{N+M}Q_j\bigr)\Bigr)
\]
使得
\[
  t(\omega) = (\omega_1,\omega_2,\dots,\omega_{N+M}),
\]
其中映射 $t$ 由 (1.6) 定义。

\noindent\textit{证明.}\quad 对元素 $f\in H^0(*\sum_{j=1}^{N+M}Q_j)$，设
$f_j(\xi_j)=\sum_n a_n^{(j)}\xi_j^n$ 是 $f$ 在点 $Q_j$ 处、关于形式参数
$\xi_j=\eta_j^{-1}(\xi)$ 的形式洛朗展开。于是
$t(f)=(f_1(\xi_1),\dots,f_{N+M}(\xi_{N+M}))$。此时有
\[
  \begin{aligned}
  \sum_{j=1}^{N+M}\operatorname{Res}_{\xi_j=0}\bigl(f_j(\xi_j)\omega_j\bigr)
  &= \sum_{j=1}^{N+M}\sum_{n\in\Z}
     (\tilde{\Psi}|\rho_j(X(n))|\tilde{u}\rangle a_n^{(j)}\\
  &= \langle\tilde{\Psi}|X\otimes t(f)|\tilde{u}\rangle = 0,
  \end{aligned}
\]
因为按我们的假设有 $\langle\Psi|X\otimes t(f)=0$。因此由定理 1.22，存在元素
$\omega\in H^0(C,\omega_C(*\sum_{j=1}^{N+M}Q_j))$ 使
$t(\omega)=(\omega_1,\dots,\omega_{N+M})$。这就证明了断言 2。\hfill$\square$
% [Tutor Note: 两处外部输入：$t$ 是 (1.6) 定义的「形式留数」映射
%  $t:\omega\mapsto(\operatorname{Res}_{Q_1}\omega,\dots)$；定理 1.22 是
%  「留数和为零的元组都能由亚纯 $1$-形式实现」的存在性定理（本质是
%  $H^1(C,\cO_C)=0$ 的 Serre 对偶形式）。第二个等号用到 $\langle\Psi|X\otimes t(f)=0$
%  —— 这是共形块的定义（$\mathfrak g\otimes\cO$ 的 $\cO$ 线性部分作用为零），
%  也是本证明唯一用到「$\Psi$ 是共形块」的地方。核对：把 $t(f)$ 取成只在
%  $Q_j$ 有非零留数的元组，则 $\langle\tilde\Psi|\rho_j(X(n))|\tilde u\rangle$
%  恰是 $\omega$ 在 $Q_j$ 的洛朗系数。原文 p.69]
%
% [Tutor Note: $\omega_j$ 左端的括号写作「$(\tilde\Psi|\ \cdots\ |\tilde u\rangle$」，
%  左端是圆括号而非角括号，与后文 $\langle\tilde\Psi|\cdots\rangle$ 不一致；
%  应为 $\langle\tilde\Psi|$。照原文排，见回报 ERRATA。原文 p.69]

\noindent\textbf{断言 3.}\quad 作为 $P$ 处（形式参数为 $\eta$）的余切向量，
$\langle\hat{\Psi}|\tilde{u}\otimes X(-1)|0\rangle d\eta$ 与断言 2 中的 $\omega$ 相合。

下面我们就把 $\omega$ 写成
\[
  \omega = \langle\tilde{\Psi}|X(P)|\tilde{u}\rangle dP .
\]

\noindent\textit{证明.}\quad 由于
$\langle\hat{\Psi}|\tilde{u}\otimes X(-1)|0\rangle d\eta$ 是 $P$ 处的余切向量，
我们可以假设 $\eta$ 是 $C$ 在 $P$ 处的局部全纯坐标。选取 $C$ 上的亚纯函数
$f\in H^0(C,\cO_C(*(P+Q_i)))$ 使得
\[
  f = \eta^{-1} + \text{在 } P \text{ 处正则},
\]
\[
  f\equiv 0 \bmod(\xi_j^{n_j}) \quad \text{在 } Q_j \text{ 处},\
  j\ne i,\ 1\le j\le N+M,
\]
其中 $n_j$ 充分大，使得 $\rho_j(X[f])|\tilde{u}\rangle=0$ 且 $f\omega$ 在 $Q_j$
处全纯，$j\ne i$，$1\le j\le N+M$。于是有
\[
  \begin{aligned}
  \langle\hat{\Psi}|\tilde{u}\otimes X(-1)|0\rangle
  &= -\sum_{k=1}^{N+M}\langle\hat{\Psi}|\rho_k(X[f])|\tilde{u}\rangle\\
  &= -\langle\tilde{\Psi}|\rho_i(X[f])|\tilde{u}\rangle .
  \end{aligned}
\]

% [原文 p.70]
另一方面，在点 $P$ 处我们有
\[
  \begin{aligned}
  \operatorname{Res}_P\Bigl(\frac{1}{\eta}\omega\Bigr)
  &= \operatorname{Res}_P(f\omega)\\
  &= -\sum_{k=1}^{N+M}\operatorname{Res}_{Q_k}(f\omega)\\
  &= -\operatorname{Res}_{Q_i}(f\omega)\\
  &= -\operatorname{Res}_{\xi_i=0}\Bigl(f_i(\xi_i)\sum_{n\in\Z}
     (\tilde{\Psi}|\rho_i(X(n))|\tilde{u}\rangle\xi_i^{-n-1}d\xi_i\Bigr)\\
  &= -\langle\tilde{\Psi}|\rho_i(X[f])|\tilde{u}\rangle\\
  &= \langle\hat{\Psi}|\tilde{u}\otimes X(-1)|0\rangle_\eta .
  \end{aligned}
\]
这就证明了断言 3。\hfill$\square$

现在我们可以证明定理 3.24 了。置
\[
  |\tilde{u}\rangle = |u\rangle\otimes X_1(-1)|0\rangle\otimes\cdots
  \otimes X_M(-1)|0\rangle .
\]
上述论证表明
\[
  \begin{aligned}
  \langle\tilde{\Psi}|\tilde{u}\rangle d\eta_1\cdots d\eta_M
  &= \langle\tilde{\Psi}|\bigl(|u\rangle\otimes X_1(-1)|0\rangle\otimes\cdots
     \otimes X_M(-1)|0\rangle\bigr) d\eta_1\cdots d\eta_M
  \end{aligned}
\]
可以看作 $T_{P_1}^*C\otimes\cdots\otimes T_{P_M}^*C$ 的元素（只要
$P_k\ne Q_j$ 且 $P_j\ne P_k$，$j\ne k$），并且关于 $P_k$ 亚纯地依赖。因此，
由 Hartogs 定理，它定义了
$H^0(C^M,\omega_C^{\boxtimes M}(\sum_{i<j}*\Delta_{ij}
+\sum_{i=1}^{M}\sum_{j=1}^{N}*\pi_i^{-1}(Q_j)))$ 中的一个元素。我们把该亚纯截面
记为
\[
  \langle\Psi|X_1(P_1)X_2(P_2)\cdots X_M(P_M)|u\rangle
  dP_1dP_2\cdots dP_M .
\]
% [Tutor Note: 这里有两步合在一起：(i) 固定 $P_1,\dots,P_{k-1},P_{k+1},\dots,P_M$，
%  上面的矩阵元是 $P_k$ 的亚纯 $1$-形式，故给出 $T^*_{P_k}C$ 的元素，于是张量起来
%  得到 $(C\setminus\{Q_j\})^M$ 上的亚纯截面；(ii) 各 $\Delta_{ij}$ 与
%  $\pi_i^{-1}(Q_j)$ 在 $C^M$ 中的余维都是 $\ge 2$（不是 $1$！），故 Hartogs 定理
%  保证截面可延拓过去。注意不能对 $C^M$ 的奇点集本身用 Hartogs（它余维 $1$），
%  原文此处实际只用到「除子外延拓」这一层。核对：$M=1$ 时无 $\Delta_{ij}$ 项，
%  得到 $H^0(C,\omega_C(*\sum_j Q_j))$，与断言 3 的 $\omega$ 一致。原文 p.70]

论断 (0) 与 (1) 由我们的定义是显然的。对论断 (2)，注意到由数据
\[
  \tilde{\mathfrak{X}} = (C;Q_1,\dots,Q_N,P_1,\dots,P_M;
  \eta_1,\dots,\eta_{N+M})
\]
定义的上述亚纯形式，与由数据
\[
  \begin{aligned}
  \tilde{\mathfrak{X}}_\sigma = (C;&Q_1,\dots,Q_N,P_{\sigma(1)},\dots,
  P_{\sigma(M)};\eta_1,\dots,\eta_N,\\
  &\eta_{N+\sigma(1)},\cdots,\eta_{N+\sigma(M)})
  \end{aligned}
\]
定义的亚纯形式相同。这就蕴涵了论断 (2)。

论断 (3) 由断言 2 推出。

下面证明论断 (4)。设点 $P'$ 落在点 $P$ 的一个小邻域 $U$ 中，$z$ 是以 $P$
为中心的局部坐标。选取亚纯函数
$f\in H^0(C,\cO_C(*(P+\sum_{k=1}^{N+M}Q_k)))$ 使得
\[
  f = z^{-1} + \text{在 } P \text{ 处正则}.
\]
此外，必要时改变 $P$ 处的局部坐标，我们可以假设 $f=z^{-1}$。于是
$w=z-z(P')$ 是 $C$ 在 $P'$ 处的局部坐标。作为
$(P,P')\times\underbrace{C\times\cdots\times C}_{M}$ 的每一点处的余切向量，
\[
  F = \langle\Psi|X(P)Y(P')X_1(P_1)\cdots X_M(P_M)|\Phi\rangle dPdP'
\]
等于
\[
  \langle\hat{\Psi}^*|(X(-1)|0_P\rangle\otimes Y(-1)|0_{P'}\rangle)
  \otimes\tilde{\Phi}\rangle dzdw
\]

% [原文 p.71]
其中
\[
  \langle\hat{\Psi}^*| = \iota(\langle\Psi|),\qquad
  \iota : \cV_\lambda^\dagger(\mathfrak{X})
  \to \cV_{\lambda,\vec{0}_{M+2}}^\dagger(\hat{\mathfrak{X}})
\]
以及
\[
  |\tilde{\Phi}\rangle = |\Phi\otimes\bigl(X_1(-1)|0_{P_1}\rangle\otimes\cdots
  \otimes X_M(-1)|0_{P_M}\rangle\bigr)
  \in\cH_{\mathfrak{X}}\otimes\cH_{\vec{0}_M} .
\]
于是我们有
\begin{equation}
  \begin{aligned}
  &\langle\hat{\Psi}^*|(X(-1)|0_P\rangle\otimes Y(-1)|0_{P'}\rangle)
   \otimes\tilde{\Phi}\rangle\\
  &\quad = -\langle\hat{\Psi}|(X[f]Y(-1)|0_{P'}\rangle)\otimes\tilde{\Phi}\rangle
    -\sum_{k=1}^{N+M}\langle\hat{\Psi}|Y(-1)|0_{P'}\rangle
    \otimes\rho_k(X[f])|\tilde{\Phi}\rangle .
  \end{aligned} \tag{3.38}
\end{equation}
% [Tutor Note: (3.38) 是把 $X[f]$ 从左边穿过 $Y(-1)|0_{P'}\rangle$ 的两种方式：
%  直接对 $\langle\hat\Psi|$ 用 Ward 恒等式得第一项（$X$ 仍留在 $P'$ 一侧），
%  或先对右边态用 $\rho_k$ 得第二项（$X$ 移到 $\tilde\Phi$ 一侧）。
%  核对方法：取 $M=0$、$\mathfrak{g}$ 阿贝尔（$[X,Y]=0$），则两边都应只剩
%  常数 $\ell(X,Y)/a^2$ 型项，这将在下面几行验证。原文 p.71]

(3.38) 的右端第二项可写成
\[
  -\sum_{k=1}^{N+M}\langle\tilde{\Psi}|Y(P')|\rho_k(X[f])\tilde{\Phi}\rangle dP'
\]
因此它在点 $P'$ 处全纯。另一方面，置 $a=z(P')$，我们有
\[
  \begin{aligned}
  (X[f])Y(-1)|0_{P'}\rangle
  &= \Bigl(X\Bigl[\frac{1}{w+a}\Bigr]\Bigr)\bigl(Y[w^{-1}]\bigr)|0_{P'}\rangle\\
  &= \Bigl(\frac{[X,Y]}{a}[w^{-1}]
     -\frac{\ell\cdot(X,Y)}{a^2}\Bigr)|0_{P'}\rangle .
  \end{aligned}
\]
因此 (3.38) 的右端第一项具有形式
\[
  \frac{\ell\cdot(X,Y)}{a^2}\langle\Psi|X_1(P_1)\cdots X_M(P_M)|\Phi\rangle
  -\frac{1}{a}\langle\Psi|[X,Y](P')X_1(P_1)\cdots X_M(P_M)|\Phi\rangle .
\]
由于 $-a=z(P)-z(P')$，我们就得到所要的结果。
类似的论证可证明论断 (5)。

\begin{remark}[3.25]
\begin{sloppypar}使用形式坐标 $\xi_j$，我们令
$(\xi_{j_1}^{\prime},\xi_{j_2}^{\prime\prime},\xi_{j_3}^{(3)},\dots,
\xi_{j_n}^{(n)})$ 是 $\underbrace{C\times\cdots\times C}_{n}$ 在中心
$(Q_{j_1},Q_{j_2},\dots,Q_{j_n})$ 处的形式坐标系，其中 $\xi_j^{(m)}$ 是
$\xi_j$ 的一份拷贝。关联函数
\end{sloppypar}
\[
  \langle\Psi|X_1(z_1)X_2(z_2)\cdots X_n(z_n)|\Phi\rangle
  dz_1dz_2\cdots dz_n
\]
在点 $(Q_{j_1},Q_{j_2},\dots,Q_{j_n})$ 处有洛朗展开
\[
  \begin{aligned}
  \sum_{m_i\in\Z}
  &\langle\Psi|\rho_1(X_1(m_1))\rho_2(X_2(m_2))\\
  &\cdots\rho_n(X_n(m_n))\Phi\rangle
  \xi_{j_1}^{\prime\,-m_1-1}\xi_{j_2}^{\prime\prime\,-m_2-1}\cdots
  \xi_{j_n}^{(n)-m_n-1}
  d\xi_{j_1}^{\prime}\xi_{j_2}^{\prime\prime}\cdots d\xi_{j_n}^{(n)}
  \end{aligned}
\]
\end{remark}
% [Tutor Note: 末行的微分只写了 $d\xi_{j_1}^{\prime}$ 与 $d\xi_{j_n}^{(n)}$，
%  中间的 $\xi_{j_2}^{\prime\prime}$ 前漏了 $d$；按张量积的定义应为
%  $d\xi_{j_1}^{\prime}d\xi_{j_2}^{\prime\prime}\cdots d\xi_{j_n}^{(n)}$。
%  见回报 ERRATA。另注意求和指标写着 $m_i\in\Z$，实指
%  $m_1,\dots,m_n\in\Z$ 的 $n$ 重和；每个因子贡献 $\xi^{-m_i-1}d\xi$，
%  故整体是 $\omega^{\boxtimes n}$ 的截面。原文 p.71]

此外，下述定理成立。

% [原文 p.72]
\begin{theorem}[3.26]
(1)\quad 置
\[
  \begin{aligned}
  &\langle\Psi|T(z)X_1(P_1)X_2(P_2)\cdots X_M(P_M)|\Phi\rangle
   dz^2dP_1dP_2\cdots dP_M\\
  &\quad = \frac{1}{2(g^*+l)}\lim_{w\to z}\Bigl\{
    \sum_{a=1}^{\dim\mathfrak{g}}
    \langle\Psi|J^a(z)J^a(w)X_1(P_1)\\
  &\quad\qquad \times X_2(P_2)\cdots X_M(P_M)|\Phi\rangle
    dzdwdP_1dP_2\cdots dP_M\\
  &\quad\qquad -\frac{\ell\dim\mathfrak{g}}{(z-w)^2}
    \langle\Psi|X_1(P_1)X_2(P_2)\cdots X_M(P_M)|\Phi\rangle
    dzdwdP_1dP_2\cdots dP_M\Bigr\}.
  \end{aligned}
\]
则这是良定义的，且对 $k=1,\dots,N$，我们有
\[
  \begin{aligned}
  &\operatorname{Res}_{\xi_k=0}\bigl(\xi_k^{n+1}
   \langle\Psi|T(\xi_k)X_1(P_1)X_2(P_2)\cdots X_M(P_M)|\Phi\rangle
   d\xi_k\bigr)dP_1dP_2\cdots dP_M\\
  &\quad = \langle\Psi|X_1(P_1)X_2(P_2)\cdots X_M(P_M)|\rho_k(L_n)\Phi\rangle
   dP_1dP_2\cdots dP_M
  \end{aligned}
\]
其中 $\{J^1,\dots,J^{\dim\mathfrak{g}}\}$ 是李代数 $\mathfrak{g}$ 的正交规范基。
（见 (2.11)。）于是我们有展开式
\[
  \begin{aligned}
  &\langle\Psi|T(\xi_k)X_1(P_1)X_2(P_2)\cdots X_M(P_M)|\Phi\rangle
   d\xi_k^2 dP_1dP_2\cdots dP_M\\
  &\quad = \sum_{n\in\Z}\langle\Psi|X_1(P_1)X_2(P_2)\cdots X_M(P_M)
   |\rho_k(L_n)\Phi\rangle\xi_k^{-n-2}
   d\xi_k^2 dP_1dP_2\cdots dP_M .
  \end{aligned}
\]

(2)\quad 对全纯坐标变换 $w=w(z)$，我们有
\[
  \begin{aligned}
  &\langle\Psi|T(w)X_1(P_1)X_2(P_2)\cdots X_M(P_M)|\Phi\rangle
   dw^2dP_1dP_2\cdots dP_M\\
  &\quad = \langle\Psi|T(z)X_1(P_1)X_2(P_2)\cdots X_M(P_M)|\Phi\rangle
   dz^2dP_1dP_2\cdots dP_M\\
  &\quad\qquad -\frac{c_v}{12}\{w(z);z\}
   \langle\Psi|X_1(P_1)X_2(P_2)\cdots X_M(P_M)|\Phi\rangle
   dz^2dP_1dP_2\cdots dP_M
  \end{aligned}
\]
其中 $\{w(z);z\}$ 是 Schwarzian 导数（Schwarzian derivative）。
（见定理 1.15。）
\end{theorem}
% [Tutor Note: 两处归一化都值得核对。(i) 系数 $1/(2(g^*+l))$ 来自 Sugawara 构造：
%  $g^*$ 是 $\mathfrak{g}$ 的对偶 Coxeter 数（$\mathfrak{g}=\mathfrak{sl}_N$ 时
%  $g^*=N$），$l$ 是水平；正交规范基指 $\sum_a (J^a,J^a)=\dim\mathfrak{g}$。
%  此时 $T(z)=\frac{1}{2(g^*+l)}\sum_a :J^aJ^a:$，其中心荷为
%  $c_v=l\dim\mathfrak{g}/(g^*+l)$。用 $\mathfrak{sl}_2$（$\dim\mathfrak{g}=3$，
%  $g^*=2$）检验：$c_v=3l/(2+l)$，与教科书一致。被减掉的
%  $l\dim\mathfrak{g}/(z-w)^2$ 即 $\sum_a l(J^a,J^a)/(z-w)^2$，是抵消
%  $J^aJ^a$ 乘积在同点处的接触奇异所必需。(ii) Schwarzian 项前的符号取决于
%  $\{w(z);z\}$ 的定义：若按通常定义
%  $\{w;z\}=w'''/w'-\frac32(w''/w')^2$（见定理 1.15），则
%  $T(w)dw^2=(T(z)-\frac{c_v}{12}\{w;z\})dz^2$，与原文一致；若某书取反号，
%  $-c_v/12$ 就会变 $+c_v/12$，这是最容易抄错的地方。数值自检：
%  $w=z+\eps$ 时 $\{w;z\}=0$；$w=\ee^z$ 时 $\{w;z\}=-1/2$，故
%  $T(\ee^z)$ 比 $T(z)$ 多出 $+c_v/24$ 的常数项。原文 p.72]

